\subsection{CT10 --- Default refinement}\label{ct:10}\label{the-default-refinement}
\ctsignature{\(P_{\ast\to10}\to \ctcert{C5};\ \ctint{P_{10\to3},P_{10\to7},P_{10\to15}};\ \ctrec{\ctnone}\).}

CT10 is an exact machine for promoting a finite missing datum to the next
branch invariant.  Its entry interface fixes the residual class, extracted
datum, finite alphabet, class table, finite-scope predicate, and C5 claim.  The
executable framework adds label-class exhaustion, finite exchange support,
direct classification, missing-datum extraction, promotion certification, and
the exact CT3, CT7, and CT15 entry contracts.

The core plan has six fields: \texttt{scope}, \texttt{labels},
\texttt{classification}, \texttt{direct}, \texttt{promotion}, and
\texttt{promotedRouting}.  The scope node produces a typed scope terminal or
a \texttt{ScopedState}.  The labels node has one successor and constructs a
\texttt{LabelState} containing both the proof that label classes exhaust the
residual class and the proof that supported exchanges are finite within each
class.

The classification node returns a C5 certificate, a
\texttt{DirectState}, or a \texttt{MissingState}.  Direct routing constructs
an exact CT3 response input or an exact CT7 class-exchange input.  A missing
state cannot route directly.  It must first pass through the one-successor
promotion node, which proves that the datum was extracted from the failed
classification and records it as a new invariant.  The promoted routing node
then constructs a provenance-tagged CT3 input or an exact CT15 rank input.

CT10 has two semantically distinct routes to the same CT3 terminal.  The
inductive \texttt{CT3Payload} records whether the input came directly from a
classified label class or only after promotion.  Each constructor retains the
corresponding predecessor state, so the consumer cannot erase this provenance.
The dependent path type permits both incoming edges while keeping every
individual execution linear.

\begin{itemize}
\tightlist
\item \textbf{S-Def} fixes the residual, promoted datum, alphabet, and class table; promotion is separately certified.
\item \textbf{S-Dich} is the exhaustive class-level classification.
\item \textbf{S-Comp} is carried by class exhaustion and the C5 certificate.
\item \textbf{S-Rout} and \textbf{S-Trig} are carried by exact CT3 and CT7 inputs and the indexed CT15 datum.
\end{itemize}

An infinite adequate refinement alphabet is the typed scope terminal.  Failure
to extract the promoted datum is not a runtime loop or unnamed residual: no
\texttt{PromotedState} can be constructed, so the tactic must be redesigned at
that node boundary.

\begin{figure}[p]
\centering
\vspace*{-1.80cm}
\resizebox{\textwidth}{!}{%
\begin{tikzpicture}[x=1cm,y=1cm,every node/.style={font=\scriptsize}]
\node[bcwide] (entry) at (0,0) {{\tiny\ttfamily CT10.entry}\\\textbf{Validated residual, datum, alphabet, table}};
\node[bcdec,bclemma,text width=3.45cm,minimum width=4.20cm] (scope) at (0,-2.30) {{\tiny\ttfamily CT10.decide.scope}\\\textbf{Finite refinement alphabet?}};
\node[bcscopefail,text width=3.10cm,minimum width=3.10cm] (scopeout) at (7.30,-2.30) {{\tiny\ttfamily CT10.terminal.scope}\\\textbf{Scope exit}};
\node[bcwide,bclemma] (labels) at (0,-4.70) {{\tiny\ttfamily CT10.certify.labels}\\\textbf{Exhaustion and finite-support certificate}};
\node[bcroutechoice,bctheorem,text width=3.30cm,minimum width=4.25cm] (classification) at (0,-7.25) {{\tiny\ttfamily CT10.decide.classification}\\\textbf{Closed, direct, or missing datum?}};
\node[bcclosed] (c5) at (-7.50,-7.25) {{\tiny\ttfamily CT10.terminal.c5}\\\textbf{C5}};
\node[bcroutechoice,bctheorem,text width=3.25cm,minimum width=4.10cm] (direct) at (-4.70,-10.20) {{\tiny\ttfamily CT10.decide.direct}\\\textbf{Route direct class result}};
\node[bcwide,bclemma,text width=3.55cm,minimum width=3.55cm] (promotion) at (4.70,-10.20) {{\tiny\ttfamily CT10.certify.promotion}\\\textbf{Promote extracted datum}};
\node[bcroutechoice,bctheorem,text width=3.35cm,minimum width=4.30cm] (promoted) at (4.70,-13.00) {{\tiny\ttfamily CT10.decide.promotedRouting}\\\textbf{Route promoted invariant}};
\node[bcroute,bcrouteneutral,bcoutputroute,text width=3.05cm,minimum width=3.05cm] (ct3) at (-4.70,-16.20) {{\tiny\ttfamily CT10.terminal.ct3}\\\textbf{CT 3 exit}};
\node[bcroute,bcrouteneutral,bcoutputroute,text width=3.05cm,minimum width=3.05cm] (ct7) at (-10.00,-13.00) {{\tiny\ttfamily CT10.terminal.ct7}\\\textbf{CT 7 exit}};
\node[bcroute,bcrouteneutral,bcoutputroute,text width=3.05cm,minimum width=3.05cm] (ct15) at (10.00,-16.20) {{\tiny\ttfamily CT10.terminal.ct15}\\\textbf{CT 15 exit}};
\draw[bcarr] (entry)--(scope);
\draw[bcscopearr] (scope)--node[bclabel,above]{exit}(scopeout);
\draw[bcarr] (scope)--node[bclabel,right]{ready}(labels);
\draw[bcarr] (labels)--node[bclabel,right]{certified}(classification);
\draw[bcarr] (classification)--node[bclabel,above]{exhausted}(c5);
\draw[bcarr] (classification)--node[bclabel,above,sloped]{direct}(direct);
\draw[bcarr] (classification)--node[bclabel,above,sloped]{missing}(promotion);
\draw[bchandoff] (direct)--node[bclabel,above,sloped]{response}(ct3);
\draw[bchandoff] (direct)--node[bclabel,above,sloped]{exchange}(ct7);
\draw[bcarr] (promotion)--node[bclabel,right]{certified}(promoted);
\draw[bchandoff] (promoted)--node[bclabel,above,sloped]{response}(ct3);
\draw[bchandoff] (promoted)--node[bclabel,above,sloped]{rank datum}(ct15);
\end{tikzpicture}%
}
\caption[Closure-tactic 10: exact refinement machine]{Closure-tactic 10 as an exact sequential machine.  Label exhaustion and datum promotion are separate one-successor certificates.  The two CT3 edges converge at one terminal but carry different constructors of \texttt{CT3Payload}, preserving direct versus promoted provenance.}
\label{fig:ct10-default-refinement}
\end{figure}
\FloatBarrier
